\documentclass[12pt]{article}

% Set landscape mode and custom margins for page while including space for the footer
\usepackage[includefoot, margin={0.5cm, 0.5cm}, landscape]{geometry}

\usepackage{graphicx}

% Used to reduce the spacing around section headers and titles.
% Format title headers by adjusting font size
% Set the spacing on all sides of the title to 0
\usepackage[explicit]{titlesec}
\titleformat{\section}{\large\bfseries}{\thesection}{1em}{\underline{#1}}
\titleformat{\subsection}{\normalsize\bfseries}{\thesubsection}{1em}{#1}
\titleformat{\subsubsection}{\small\bfseries}{\thesubsubsection}{1em}{#1}
\titlespacing*{\section}{0pt}{*0}{0pt}
\titlespacing*{\subsection}{0pt}{*0}{0pt}
\titlespacing*{\subsubsection}{0pt}{*0}{0pt}

% Reduce spacing around list items
\usepackage{enumitem}
\setlist{nolistsep}

% Set line spacing
\usepackage{setspace}
\singlespacing

% Create a multicolumn layout
% Set the amount of separation between columns
% Draw a vertical rule between the columns
\usepackage{multicol}
\setlength{\columnsep}{1cm}
\setlength{\columnseprule}{0.1pt}

% Add a frame around the content
\usepackage{mdframed}

% Remove paragraph indent
\setlength{\parindent}{0pt}

% Custom footer (and header if I wanted)
\usepackage{fancyhdr}
\pagestyle{fancy}

% Set custom headers and footers for fancyhdr
\fancyhead{}
\fancyfoot[C]{Primal: The Awakening Summary Sheet v1.0}
\renewcommand{\headrulewidth}{0pt} 	% Remove horizontal rule from header
\renewcommand{\footrulewidth}{0pt}  % Remove horizontal rule from footer

% Custom frame style for mdframed
% The negative margin is needed to fix a weird spacing that I couldn't figure out
\mdfdefinestyle{customFrame}{%
    outerlinewidth = 0.4pt,
    innertopmargin = -0.3cm}

\mdfdefinestyle{SummaryCard}{
    linecolor=black,
    outerlinewidth=2pt,
    roundcorner=20pt,
    innertopmargin=4pt,
    innerbottommargin=4pt,
    innerrightmargin=4pt,
    innerleftmargin=4pt,
    leftmargin = 4pt,
    rightmargin = 4pt,
    frametitlealignment=\center
}

% Reduce whitespace in the enumerate list environment
\newenvironment{enumerateCustom}
{\begin{enumerate}
  \setlength{\itemsep}{1pt}
  \setlength{\parskip}{0pt}
  \setlength{\parsep}{0pt}}
{\end{enumerate}}

% Reduce whitespace in the itemize list environment
\newenvironment{itemizeCustom}
{\begin{itemize}
  \setlength{\itemsep}{1pt}
  \setlength{\parskip}{0pt}
  \setlength{\parsep}{0pt}}
{\end{itemize}}

\begin{document}

\begin{mdframed}[style=SummaryCard, align=center, userdefinedwidth=35em, frametitle={Round Order Guide}]

    \begin{enumerateCustom}
        \item Each player may resolve one \textbf{Consume} ability.
        \item Monster Upkeep
            \begin{itemizeCustom}
                \item Replace behavior card with lowest number (or all tied cards).
                \item Gain 1 per player (1\includegraphics[scale=0.40]{images/per_player.png}) \textbf{struggle} tokens.
                    \begin{itemizeCustom}
                        \item If $\ge$ 3\includegraphics[scale=0.40]{images/per_player.png}, \textbf{Unleash}
                    \end{itemizeCustom}
            \end{itemizeCustom}

        \item Player Turns: Clockwise turn order starting with \textbf{aggro} token player.
            \begin{itemizeCustom}
                \item Movement Phase
                    \begin{itemizeCustom}
                        \item May spend 1 \textbf{stamina} to move to adjacent sector.
                            \begin{itemizeCustom}
                                \item \textit{Remember, remove threatened status when you move.}
                            \end{itemizeCustom}
                        \item If don't move: take \textbf{threatened} token.
                    \end{itemizeCustom}
                \item Action Phase: Perform actions in any order and count.
                    \begin{itemizeCustom}
                        \item Play action card.
                            \begin{itemizeCustom}
                                \item Play up to 5 action cards into your \textbf {sequence}. Resolve card ability text and then specific card color effect.
                                    \begin{itemizeCustom}
                                        \item Take \textbf{aggro} token immediately when played, if relevant.
                                        \item Player chooses order to resolve simultaneous triggered abilities on played cards.
                                    \end{itemizeCustom}
                            \end{itemizeCustom}
                        \item Activate \textbf{Action} ability
                        \item Revive (limit once per turn)

                        \textbf{End of Phase:} Gain \textbf{stamina} token if $\ge 2$ cards in hand.
                    \end{itemizeCustom}
                \item Attrition Phase
                    \begin{itemizeCustom}
                        \item Reveal 1 attrition card (2, pick highest if \textbf{threatened})
                        \item Take \textbf{attrition damage} if number of defense cards in sequence $<$ value. See damage value on bottom-right of stance card.
                    \end{itemizeCustom}
                \item End of Turn
                    \begin{itemizeCustom}
                        \item Discard sequence in order played
                        \item Draw / discard to hand limit (default 5)
                            \begin{itemizeCustom}
                                \item When deck is empty, suffer damage equal to weapon level and then reshuffle.
                            \end{itemizeCustom}
                        \item Rotate monster to face \textbf{aggro} player.
                    \end{itemizeCustom}
            \end{itemizeCustom}

        \item End of Round
            \begin{itemizeCustom}
                \item Resolve "at the end of the round" effects
                \item Advance round marker.
            \end{itemizeCustom}
    \end{enumerateCustom}
\end{mdframed}

\pagebreak

\begin{mdframed}[style=SummaryCard, align=center, userdefinedwidth=35em, frametitle={Important Keywords (Mirah)}]
    \begin{description}
        \item[Assist] During another player’s turn, either before or after any action, you may discard an Assist card. The active player draws 1. Limit once per turn.
        \item[Blind] When you inflict Blind, choose a peril card in play. That peril is considered to be blank until the start of your next turn.
        \item[Focused Mastery] \includegraphics[scale=0.30]{images/mastery_icon.png} If card ability text begins with this icon, it means you can only resolve the ability if your mastery card is focused.
        \item[Recycle] When you are instructed to Recycle X, you may discard X cards of your choice from your hand to draw X.
        \item[Stealth] When you play an action card with the keyword Stealth in your sequence, that card do not trigger reaction icons containing color or card type icons.
        \item[Volley] You may discard X cards from the top of your deck and deal damage to the monster equal to your weapon level for each offensive card (red / blue) discarded this way.
        \item[Weapon Icon] \includegraphics[scale=0.20]{images/weapon_level.png} These are the class icons that represent the different hunter classes. They refer to your weapon level (top right corner of weapon card).
    \end{description}
\end{mdframed}

\pagebreak

\begin{mdframed}[style=SummaryCard, align=center, userdefinedwidth=35em, frametitle={Important Keywords (Dareon)}]
    \begin{description}
        \item[Assist] During another player’s turn, either before or after any action, you may discard an Assist card. The active player draws 1. Limit once per turn.
        \item[Berserker] This is a keyword ability. When you play a card with the keyword Berserker in your sequence, you enter the berserker state and remain in that state until the end of your turn. While in the berserker state, draw 1 each time you play an attack card in your sequence.
        \item[Finisher] After you play a card with the keyword Finisher in your sequence, your Action phase ends
        \item[Focused Mastery] \includegraphics[scale=0.30]{images/mastery_icon.png} If card ability text begins with this icon, it means you can only resolve the ability if your mastery card is focused.
        \item[Resilience] After you suffer , you may reveal a card with the keyword Resilience from your hand to draw 1. Limit once per Resilience card per round.
        \item[Vulnerable] Inflict Vulnerable on the monster. While the monster is vulnerable, you may discard the vulnerable token to double a single source of damage. Vulnerability is removed at the end of the round.
        \item[Weapon Icon] \includegraphics[scale=0.20]{images/weapon_level.png} These are the class icons that represent the different hunter classes. They refer to your weapon level (top right corner of weapon card).
    \end{description}
\end{mdframed}

\pagebreak

\begin{mdframed}[style=SummaryCard, align=center, userdefinedwidth=35em, frametitle={Important Keywords (Thoreg)}]
    \begin{description}
        \item[Assist] During another player’s turn, either before or after any action, you may discard an Assist card. The active player draws 1. Limit once per turn.
        \item[Confuse] Give monster the confused token and turn monster to a sector of your choice. When the monster would activate a boost effect, discard the confused token instead and cancel that effect (still pay the boost cost). Remove at beginning of next round.
        \item[Focused Mastery] \includegraphics[scale=0.30]{images/mastery_icon.png} If card ability text begins with this icon, it means you can only resolve the ability if your mastery card is focused.
        \item[Recycle] When you are instructed to Recycle X, you may discard X cards of your choice from your hand to draw X.
        \item[Stealth] When you play an action card with the keyword Stealth in your sequence, that card do not trigger reaction icons containing color or card type icons.
        \item[Strain] Then the next time you refill your hand, draw one card fewer for each strain you have.
        \item[Stun] When you stun the monster, place stun token on an active behavior card. That card cannot be triggered this round. Then choose a player. They reveal up to two attack cards from their hand and deal weapon damage per revealed card. Remove after resolving "end of round" effects.
        \item[Weapon Icon] \includegraphics[scale=0.20]{images/weapon_level.png} These are the class icons that represent the different hunter classes. They refer to your weapon level (top right corner of weapon card).
    \end{description}
\end{mdframed}

\pagebreak

\begin{mdframed}[style=SummaryCard, align=center, userdefinedwidth=35em, frametitle={Important Keywords (Ljonar)}]
    \begin{description}
        \item[Assist] During another player’s turn, either before or after any action, you may discard an Assist card. The active player draws 1. Limit once per turn.
        \item[Confuse] Give monster the confused token and turn monster to a sector of your choice. When the monster would activate a boost effect, discard the confused token instead and cancel that effect (still pay the boost cost). Remove at beginning of next round.
        \item[Focused Mastery] \includegraphics[scale=0.30]{images/mastery_icon.png} If card ability text begins with this icon, it means you can only resolve the ability if your mastery card is focused.
        \item[Resilience] After you suffer \includegraphics[scale=0.30]{images/monster_damage.png}, you may reveal a card with the keyword Resilience from your hand to draw 1. Limit once per Resilience card per round.
        \item[Taunt] During another player’s turn, discard a Taunt card. Active player draws 1, you take the aggro token, and turn monster to your sector. Limit once per turn.
        \item[Vulnerable] Inflict Vulnerable on the monster. While the monster is vulnerable, you may discard the vulnerable token to double a single source of damage. Vulnerability is removed at the end of the round.
        \item[Weapon Icon] \includegraphics[scale=0.20]{images/weapon_level.png} These are the class icons that represent the different hunter classes. They refer to your weapon level (top right corner of weapon card).
    \end{description}
\end{mdframed}

\pagebreak

\begin{mdframed}[style=SummaryCard, align=center, userdefinedwidth=35em, frametitle={Important Triggers and Resolutions}]
%\begin{multicols*}{2}

%\section*{Important Triggers and Resolutions}

    \begin{description}
        \item[Empty Player Deck] Suffer damage equal to your weapon level. Immediately reshuffle your deck.
        \item[Empty Behavior Deck] Immediately reshuffle discard pile. Monster gains 1 struggle.
        \item[Unleashed] If struggle $\ge$ 3\includegraphics[scale=0.40]{images/per_player.png}:
            \begin{enumerateCustom}
                \item All players suffer \includegraphics[scale=0.30]{images/monster_damage.png}
                \item Remove all but 1\includegraphics[scale=0.40]{images/per_player.png} struggle tokens.
            \end{enumerateCustom}
        \item[Card Resolution Timing] Resolve simultaneous triggers with the following priority. If multiple cards have the same priority, the players decide the resolution order.
            \begin{enumerateCustom}
                \item Stance Cards
                \item Peril Cards
                \item Behavior Cards
                \item Other Cards
            \end{enumerateCustom}
    \end{description}
%\section*{Goal of the Game}
%\section*{Player Turn}
%\section*{End Game}

%\end{multicols*}
\end{mdframed}
\end{document}
